% © 2026 Ing. David Jaroš — CC BY-NC-ND 4.0
%
% This work is licensed under a Creative Commons Attribution-NonCommercial-NoDerivatives
% 4.0 International License (CC BY-NC-ND 4.0).
%
% R ≈ 1.114 — New Approaches (Open Problem B): F1–F4
% Track: CORE — α derivation — v61 baseline
% Date: 2026-03-08

\ifdefined\INCLUDEMODE
  % Being included in another document — skip preamble
\else
  \documentclass[11pt,a4paper]{article}
  \usepackage{amsmath,amssymb,amsthm}
  \usepackage{hyperref}
  \usepackage{booktabs}

  \newtheorem{theorem}{Theorem}
  \newtheorem{proposition}{Proposition}
  \newtheorem{conjecture}{Conjecture}
  \theoremstyle{remark}
  \newtheorem{remark}{Remark}
  \newtheorem{gap}{Gap}

  \title{$R \approx 1.114$: New Approaches (Open Problem B)\\
    \large Approaches F1 (Two-Loop KK), F2 (Modular Weights),\\
    F3 (Ramanujan $\tau$), and F4 ($\psi$-Circle Holonomy)}
  \author{Ing.\ David Jaroš}
  \date{2026-03-08}

  \begin{document}
  \maketitle
\fi

% ─────────────────────────────────────────────────────────────────────────────
\section{Problem Statement}
\label{sec:r_new_problem}

The fine-structure constant derivation in UBT requires
\begin{equation}
  B = B_{\mathrm{base}} \times R
    = N_{\mathrm{eff}}^{3/2} \times R
    \approx 41.57 \times 1.114
    \approx 46.3.
\end{equation}
The correction factor $R \approx 1.114$ is phenomenological (\textbf{OPEN PROBLEM B}).
Previous approaches (B1: volume ratios, B2: algebraic candidates, B3: two-loop
on $\psi$-circle) have been documented in \texttt{r\_factor\_geometry.tex}.
Best known candidate: $R \approx 1 + \alpha(N_{\mathrm{eff}} + \pi + 1/4) \approx
1.11235$, error $0.15\%$, status \textbf{[NUMERICAL OBSERVATION]}.

This document explores four new approaches F1–F4, all avoiding $n^* = 137$
or $\alpha = 1/137$ as inputs (non-circularity requirement).

\begin{description}
  \item[F1] $R$ from two-loop QED correction on the compact $\psi$-circle
            with discrete Kaluza--Klein spectrum.
  \item[F2] $R$ as a ratio of modular weights from the Hecke analysis.
  \item[F3] $R$ from the Ramanujan $\tau$-function.
  \item[F4] $R$ from $\psi$-circle holonomy (Hosotani phase).
\end{description}

% ─────────────────────────────────────────────────────────────────────────────
\section{Approach F1: Two-Loop QED Correction with KK Spectrum}
\label{sec:F1_twoloop_KK}

\textbf{Status: [DEAD END]} — the required two-loop coefficient $g(1) \approx 49$
does not emerge naturally from known QED two-loop structures.

\subsection{Setup}

The imaginary-time circle $S^1_\psi$ of radius $R_\psi = \hbar/(m_e c)$
quantises the Kaluza--Klein (KK) spectrum.  The standard two-loop vacuum
polarisation acquires a finite-size correction:
\begin{equation}
  B = B_{1\text{-loop}} \times \Bigl[1 + \frac{\alpha}{\pi}\,f(m_e R_\psi)\Bigr].
\end{equation}
In the $\psi$-circle compactification, the loop integral splits as
\begin{equation}
  \int\!\frac{d^4 k}{(2\pi)^4} \;\longrightarrow\;
  \frac{1}{2\pi R_\psi}\sum_{n=-\infty}^{\infty}\int\!\frac{d^3 k}{(2\pi)^3},
  \qquad k_\psi = \frac{n}{R_\psi}.
\end{equation}
For $R_\psi = \hbar/(m_e c)$, we have $m_e R_\psi = 1$ in natural units, so:
\begin{equation}
  R = 1 + \frac{\alpha}{\pi}\,g(1),
  \qquad
  g(1) = \frac{(R-1)\,\pi}{\alpha}
        = \frac{0.114 \times \pi}{\alpha}
        \approx 0.114 \times 430 \approx 49.
  \label{eq:F1_g1_required}
\end{equation}

\subsection{Numerical Exploration}

The standard two-loop QED corrections to the photon propagator give
coefficients such as
\begin{equation}
  C_{2\text{L}} = \frac{11}{12}\,N_{\mathrm{eff}} \approx 11 \quad
  \text{(for } N_{\mathrm{eff}} = 12\text{)},
\end{equation}
which falls short of the required $g(1) \approx 49$ by a factor of $\sim 4.5$.

\begin{center}
\begin{tabular}{lll}
  \toprule
  $g(1)$ candidate & Value & Error vs.\ required $49$ \\
  \midrule
  $\frac{11}{12} N_{\mathrm{eff}}$          & $11$      & $-78\%$ \\
  $4\pi$                                     & $12.57$   & $-74\%$ \\
  $N_{\mathrm{eff}}^2/4$                    & $36$      & $-27\%$ \\
  $N_{\mathrm{eff}}^2/3$                    & $48$      & $-2\%$  \\
  $N_{\mathrm{eff}}^2/(3 - \epsilon)$       & depends   & open    \\
  \bottomrule
\end{tabular}
\end{center}

The closest algebraic candidate is $g(1) \approx N_{\mathrm{eff}}^2/3 = 48$
(error $2\%$), which would give $R = 1 + \alpha N_{\mathrm{eff}}^2/(3\pi)
\approx 1.111$ (error $0.27\%$).  However, the combination
$N_{\mathrm{eff}}^2/3$ lacks physical motivation in the two-loop QED context.

\subsection{Conclusion for F1}

$g(1) \approx 49$ is \emph{not} a natural two-loop QED coefficient.  The
one closest standard result ($\frac{11}{12}N_{\mathrm{eff}} = 11$) is off by
a factor $\sim 4.5$, and no combination of known two-loop coefficients with
$N_{\mathrm{eff}} = 12$ recovers $g(1) \approx 49$ from first principles.
This approach is classified as \textbf{[DEAD END]}.

% ─────────────────────────────────────────────────────────────────────────────
\section{Approach F2: R as a Ratio of Modular Weights}
\label{sec:F2_modular_weights}

\textbf{Status: [DEAD END]} — no natural ratio of the weights $\{2, 4, 6\}$
and $N_{\mathrm{eff}} = 12$ gives a value in $[1.09, 1.14]$.

\subsection{Setup}

Both Set A and Set B Hecke forms (UBT lepton conjecture) have weights
$k = 2, 4, 6$, whose sum equals $12 = N_{\mathrm{eff}}$.  A systematic
search for a ratio built from $\{2, 4, 6\}$ and $N_{\mathrm{eff}} = 12$
that yields $R \approx 1.114$ was performed.

\subsection{Systematic Search}

\begin{center}
\begin{tabular}{llll}
  \toprule
  Formula & Value & Error vs.\ $1.114$ & Status \\
  \midrule
  $k_{\max}/k_{\mathrm{avg}} = 6/4$           & $1.500$  & $+35\%$ & FAIL \\
  $(k_1+k_2+k_3)/(3\,k_1) = 12/6$            & $2.000$  & $+79\%$ & FAIL \\
  $\sqrt{k_{\mathrm{avg}}/k_{\min}} = \sqrt{4/2}$ & $1.414$ & $+27\%$ & FAIL \\
  $1 + k_{\min}/k_{\max} = 1 + 2/6$          & $1.333$  & $+20\%$ & FAIL \\
  $k_{\mathrm{avg}}/k_{\min} \cdot 1/4 = 4/2/4$ & $0.500$ & $-55\%$ & FAIL \\
  $1 + 1/N_{\mathrm{eff}} = 1 + 1/12$        & $1.083$  & $-2.8\%$ & [CLOSE] \\
  $1 + k_{\min}/N_{\mathrm{eff}} = 1 + 2/12$ & $1.167$  & $+4.8\%$ & FAIL \\
  $(k_{\max} + k_{\min})/(k_{\mathrm{avg}}+k_{\min}) = 8/6$ & $1.333$ & $+20\%$ & FAIL \\
  $N_{\mathrm{eff}}^{1/2}/k_{\mathrm{avg}} = \sqrt{12}/4$ & $0.866$ & $-22\%$ & FAIL \\
  \bottomrule
\end{tabular}
\end{center}

The closest result is $1 + 1/N_{\mathrm{eff}} = 13/12 \approx 1.083$
(error $-2.8\%$), which is outside the target window $[1.09, 1.14]$.

\begin{remark}
  The value $1 + 1/N_{\mathrm{eff}} = 13/12$ is tantalisingly close and
  has a natural interpretation: a fractional KK mode $1/N_{\mathrm{eff}}$
  on the $\psi$-circle.  However, the $2.8\%$ deviation cannot be dismissed
  as a rounding error, and no derivation of this formula from
  $\mathcal{S}[\Theta]$ exists.  This is a \textbf{[NUMERICAL OBSERVATION]},
  not a proof.
\end{remark}

\subsection{Conclusion for F2}

No natural ratio of the Hecke weights $\{2,4,6\}$ and $N_{\mathrm{eff}} = 12$
reproduces $R = 1.114$ within $1\%$.  This approach is classified as
\textbf{[DEAD END]}.

% ─────────────────────────────────────────────────────────────────────────────
\section{Approach F3: R from the Ramanujan $\tau$-Function}
\label{sec:F3_ramanujan}

\textbf{Status: [DEAD END]} — no combination of $\tau(p)$ values and
$N_{\mathrm{eff}}$ recovers $R = 1.114$ without using $\alpha$ as input.

\subsection{Background}

The Ramanujan $\tau$-function gives the Fourier coefficients of the unique
weight-12 cusp form $\Delta(z) = q\prod_{n=1}^\infty(1-q^n)^{24}$:
\begin{equation}
  \tau(2) = -24, \quad \tau(3) = 252, \quad \tau(5) = 4830,
  \quad \tau(7) = -16744, \quad \tau(11) = 534612.
\end{equation}
These are intimately related to $N_{\mathrm{eff}} = 12 = k_{\Delta}$
(the weight of $\Delta$) and the number $24 = |\tau(2)|$.

\subsection{Systematic Search}

\begin{center}
\begin{tabular}{llll}
  \toprule
  Formula & Value & Error vs.\ $1.114$ & Status \\
  \midrule
  $1 + |\tau(2)|/N_{\mathrm{eff}}^2 = 1 + 24/144$ & $1.167$ & $+4.8\%$ & FAIL \\
  $1 + |\tau(2)|/(N_{\mathrm{eff}} \cdot 17) = 1 + 24/204$ & $1.118$ & $+0.3\%$ & [CLOSE] \\
  $1 + |\tau(2)|/(N_{\mathrm{eff}} \cdot 18)$       & $1.111$ & $-0.3\%$ & [CLOSE] \\
  $1 + |\tau(2)|/(N_{\mathrm{eff}} \cdot k_\Delta)$ & $1 + 24/144 = 1.167$ & $+4.8\%$ & FAIL \\
  $|\tau(2)|/(N_{\mathrm{eff}} + |\tau(2)|) = 24/36$ & $0.667$ & $-40\%$ & FAIL \\
  $\tau(3)/\tau(2)^2 = 252/576$                      & $0.438$ & $-61\%$ & FAIL \\
  \bottomrule
\end{tabular}
\end{center}

Two close observations:
\begin{align}
  1 + \frac{|\tau(2)|}{N_{\mathrm{eff}} \times 17} &= 1 + \frac{24}{204} \approx 1.118, \quad \text{error } +0.3\%, \\
  1 + \frac{|\tau(2)|}{N_{\mathrm{eff}} \times 18} &= 1 + \frac{24}{216} \approx 1.111, \quad \text{error } -0.3\%.
\end{align}
Both are \textbf{[NUMERICAL OBSERVATIONS]}.  The denominator factor 17
(or 18) has no clear physical motivation.

\begin{remark}
  $|\tau(2)| = 24$ is also the dimension of the Leech lattice \(\Lambda_{24}\),
  connected to the Monster group and moonshine.  The appearance of 24 in the
  denominator of a formula close to $R$ is intriguing but purely numerical.
  No algebraic derivation from $\mathcal{S}[\Theta]$ exists at this time.
\end{remark}

\subsection{Conclusion for F3}

The Ramanujan $\tau$-function provides no derivation of $R = 1.114$ from
first principles without using $\alpha$ or other fitted inputs.  This approach
is classified as \textbf{[DEAD END]}, though the observation
$1 + 24/(N_{\mathrm{eff}} \times 17) \approx 1.118$ at $0.3\%$ error is
recorded as a \textbf{[NUMERICAL OBSERVATION]}.

% ─────────────────────────────────────────────────────────────────────────────
\section{Approach F4: R from $\psi$-Circle Holonomy}
\label{sec:F4_holonomy}

\textbf{Status: [DEAD END]} — $|2|^x = 1.114$ requires $x = \ln(1.114)/\ln 2
\approx 0.155$, which is not a simple fraction.

\subsection{Setup}

On the compact $\psi$-circle the gauge holonomy is
\begin{equation}
  U = \exp\!\Bigl(i\oint A_\psi\,d\psi\Bigr) = \exp(2\pi i\,\theta_H),
\end{equation}
where $\theta_H$ is the Hosotani phase.  For a $U(1)$ gauge field with
an electron ($Q = -1$) in the fundamental representation,
$\theta_H = 1/2$ (mod 1), giving $U = \exp(i\pi) = -1$.

\subsection{Holonomy-Based Candidates}

\begin{center}
\begin{tabular}{llll}
  \toprule
  Formula & Value & Error vs.\ $1.114$ & Status \\
  \midrule
  $|1 - \exp(i\pi)| = |1 - (-1)|$         & $2$       & $+79\%$ & FAIL \\
  $|1 + \exp(i\pi)| = |1 - 1|$            & $0$       & $-100\%$ & FAIL \\
  $|1 - \exp(i\pi)|^{1/6} = 2^{1/6}$      & $1.122$   & $+0.76\%$ & [CLOSE, known B2] \\
  $|1 - \exp(i\pi)|^x$, $x = 0.155$       & $1.114$   & exact fit & no simple $x$ \\
  $|\det(1+U)|^{1/N_{\mathrm{eff}}}$ ($U$ = $2\times 2$ matrix $-I$)
    & $|0|^{1/12} = 0$ & $-100\%$ & FAIL \\
  $|\mathrm{tr}(1 + U)|^{1/N_{\mathrm{eff}}} = |0|^{1/12}$
    & $0$ & $-100\%$ & FAIL \\
  \bottomrule
\end{tabular}
\end{center}

The only viable formula is $2^{1/6}$ (known approach B2, documented in
\texttt{r\_factor\_geometry.tex}), which requires $x = 1/6$.  For $R =
1.114$ exactly one needs $x = \ln(1.114)/\ln 2 \approx 0.1549$; the
nearest simple fractions are:
\begin{equation}
  x = \frac{1}{6} \approx 0.1667 \;\text{(error }7.6\%\text{ on }x\text{)},
  \qquad
  x = \frac{3}{19} \approx 0.1579 \;\text{(error }1.9\%\text{ on }x\text{)}.
\end{equation}
Neither $1/6$ nor $3/19$ has been derived from $\mathcal{S}[\Theta]$.

\subsection{Conclusion for F4}

Holonomy-based formulas either reproduce the known $2^{1/6}$ candidate
(B2) or give trivial/zero values.  The required exponent $x \approx 0.155$
is not a simple fraction.  This approach is classified as
\textbf{[DEAD END]} (with $2^{1/6}$ already documented in B2).

% ─────────────────────────────────────────────────────────────────────────────
\section{Summary}
\label{sec:r_new_summary}

\begin{center}
\begin{tabular}{lll}
  \toprule
  Approach & Best result & Status \\
  \midrule
  F1: Two-loop KK on $\psi$-circle
    & $g(1) \approx N_{\mathrm{eff}}^2/3 = 48$, $R \approx 1.111$, err $0.27\%$
    & [DEAD END] \\
  F2: Modular weight ratios
    & $1 + 1/N_{\mathrm{eff}} = 13/12 \approx 1.083$, err $-2.8\%$
    & [DEAD END] \\
  F3: Ramanujan $\tau$-function
    & $1 + 24/(N_{\mathrm{eff}} \times 17) \approx 1.118$, err $+0.3\%$
    & [NUMERICAL OBS, DEAD END] \\
  F4: $\psi$-circle holonomy
    & $2^{1/6} \approx 1.122$, err $0.76\%$ (= known B2)
    & [DEAD END] \\
  \midrule
  B3 best (from \texttt{r\_factor\_geometry.tex})
    & $1 + \alpha(N_{\mathrm{eff}} + \pi + 1/4) \approx 1.11235$, err $0.15\%$
    & [NUMERICAL OBS] \\
  \bottomrule
\end{tabular}
\end{center}

\textbf{Open Problem B remains fully open after eight approaches (B1, B2, B3,
F1, F2, F3, F4 and the algebraic scan).}  No approach has produced a derivation
of $R = 1.114$ from $\mathcal{S}[\Theta]$ without using $\alpha$ as an input.

\begin{gap}
  The geometric or algebraic origin of $R \approx 1.114$ is unknown.
  The most promising numerical observations are:
  \begin{enumerate}
    \item $1 + \alpha(N_{\mathrm{eff}} + \pi + 1/4) \approx 1.11235$ (err $0.15\%$,
          requires $\alpha$ as input — circular).
    \item $1 + 24/(N_{\mathrm{eff}} \times 17) \approx 1.118$ (err $0.3\%$,
          denominator factor 17 unmotivated).
    \item $1 + 1/N_{\mathrm{eff}} = 13/12 \approx 1.083$ (err $-2.8\%$,
          simplest formula but outside the target window).
  \end{enumerate}
  A derivation from first principles requires identifying a geometric or
  algebraic quantity in $\mathcal{S}[\Theta]$ that evaluates to $\approx 0.114$
  without using $n^* = 137$ or $\alpha$ as inputs.
\end{gap}

\ifdefined\INCLUDEMODE\else
\end{document}
\fi
